\chapter{Multiple Integrals, Change of Variables, and Vector Field Theorems}
\chaptermark{Multiple Integrals and Vector Field Theorems}

A cooling plate is a disk of radius one meter. Its thickness is
uniform, but its density falls off from the center: $e^{-(x^{2}+y^{2})}$ kilograms
per square meter. Its mass is the sum of the density over the whole disk: a numerical library reports
$1.9858653038$ kilograms after $11907$ function evaluations.

A colleague writes a different version of the same computation. He moves to
polar coordinates and writes
\[
\int_{0}^{2\pi}\!\!\int_{0}^{1} e^{-r^{2}}\,dr\,d\theta ,
\]
and gets $4.6924344183$ kilograms. His program runs without a warning.
That number is $2.36$ times larger than the first, and not one of the
programs complains. The mistake is one letter: the factor $r$ is missing, so the area
of every small cell is counted too large.

Damage of that size also shows up on the other side of this chapter. Compute the work of
$F(x,y)=(-y,x)/(x^{2}+y^{2})$ around the unit circle: the result is
$6.2831853072$, even though the curl of that field is zero at every point where it
is defined. If zero curl meant zero closed work, one of those two computations
would lie. What lies is not the computation but the hypothesis that was left
unchecked: the field is not defined at the origin, and the origin lies inside
the circle.

Chapter 13 built the integral on an interval, and Chapter 12 gave the Jacobian as the
matrix of derivatives. Neither is enough yet: there is still no way to add up over a region
of dimension greater than one, no rule for what happens to the area of a cell
when the coordinates are changed, and no meaning for the sign on a boundary. This chapter
fills in all three, then shows that Green, Gauss, and Stokes are in fact one
statement.

\begin{untukanak}
Picture a tailor with a shirt pattern on squared paper. The pattern is square
and regular; the cloth to be cut curves to follow the body. If that square pattern
is pasted on as it is, the shirt is tight on one side and loose on the other,
even though the number of cells stays the same.

An experienced tailor does something else: every small cell will stretch
or shrink by a certain ratio as it is carried onto the cloth. Cells near the waist stretch
a little; cells near the shoulder shrink. He records that ratio, multiplies the area
of the cell by it, and only then adds up.

The same work appears with curved shapes: many objects, areas, or weights.
If the stretch ratio is forgotten, the total is still a tidy number, and
precisely because it is tidy it is hard to distrust. In this chapter we shall learn why
that ratio has to be there, and how to compute it for round shapes.
\end{untukanak}

\section{From rectangles to general regions}
\tingkat{basic}

We begin with the simplest object: a rectangle, and a
number we want to add up over its whole surface. The method is the same as in
Chapter 13, only the pieces are now cells rather than segments.

\begin{definisi}[The double integral]\index{double integral}
Let $R=[a,b]\times[c,d]$ be a closed rectangle. A
\emph{partition} $\mathcal{P}$ of $R$ is a division of $[a,b]$ into
intervals and of $[c,d]$ into intervals; the result is $mn$ small
rectangles $R_{ij}$. The \emph{mesh} of the partition is the longest diagonal among
those cells, and is written $\|\mathcal{P}\|$.

For $f$ bounded on $R$, write $M_{ij}=\sup_{R_{ij}}f$ and
$m_{ij}=\inf_{R_{ij}}f$, then
\[
U(\mathcal{P},f)=\sum_{i,j}M_{ij}\,\mathrm{luas}(R_{ij}),
\qquad
L(\mathcal{P},f)=\sum_{i,j}m_{ij}\,\mathrm{luas}(R_{ij}).
\]
The function $f$ is called \emph{integrable} on $R$ when the supremum of all $L$ equals
the infimum of all $U$; that common value is $\iint_{R}f\,dA$.
\end{definisi}

This definition says nothing about the order of integration. If $f$ is continuous, the upper
and lower sums meet because the largest gap of $f$ on a single cell
shrinks with the mesh: on a compact set, a continuous function is uniformly continuous
(Chapter 11).

The next question is practical. To add up over cells means computing
something two-dimensional, whereas our hands can only compute one-variable
integrals. The bridge: slice along one direction first, then add up
the results of those slices.

\begin{teorema}[Fubini]\index{Fubini's theorem}
If $f$ is continuous on $R=[a,b]\times[c,d]$, then
\[
\iint_{R} f\,dA=\int_{a}^{b}\!\!\int_{c}^{d} f(x,y)\,dy\,dx
=\int_{c}^{d}\!\!\int_{a}^{b} f(x,y)\,dx\,dy .
\]
\end{teorema}

\begin{proof}
A full proof demands matching the two-dimensional upper and lower sums with the one-dimensional
integral uniformly, and it is in Spivak, the chapter \emph{Integration},
Theorem 3-10. The idea runs like this.

For $f\ge 0$, the double integral is the volume of the solid under the graph of $f$.
Slice that solid with a plane of constant $x$. The slice is a region under the curve
$y\mapsto f(x,y)$, so its area is
\[
A(x)=\int_{c}^{d}f(x,y)\,dy .
\]
The volume of the solid is the integral of the area of the slice:
\[
\iint_{R}f\,dA=\int_{a}^{b}A(x)\,dx=\int_{a}^{b}\!\!\int_{c}^{d}f(x,y)\,dy\,dx .
\]
Because $f$ is continuous on a compact set, $A$ is continuous. For $f$ of any
sign, split $f=f^{+}-f^{-}$; the other direction comes from slices of constant $y$.
\end{proof}

What matters in this theorem is not the formula but the permission it grants: the order may
be swapped. That permission lapses for functions that are not integrable, and that is
discussed later.

The solid sliced in that proof appears in Figure~\ref{fig:kolom-tipis}:
each column contributes $f(x,y)\,\Delta A$.

\begin{bacakeras}
--- How much sand is in this pile?
--- Add up its thin columns. Stand on the north side, slice the pile into
strips, measure each strip, then add.
--- What do you get?
--- Eighteen buckets.
--- Now stand on the east side, slice in the other direction, measure, add.
--- What do you get?
--- Eighteen buckets as well.
--- The north-south strips and the east-west strips are not the same at all. Why
is the answer the same?
--- Because the pile is one. We are not adding two piles; we are slicing
the same object with two knives in different directions.
--- So the direction of the slice may be swapped?
--- As long as the object is measurable. If the pile has a hole in the middle, the two
directions still agree; what fails later is the hole, not the order.

Say it once more, in an ordinary voice: the direction of the slice changes, the volume stays.
\end{bacakeras}

\begin{figure}[htbp]
\centering
\begin{tikzpicture}[scale=1.55,font=\small]
  \fill[layergrey!10] (0,0) -- (6.1,0) -- (7.1,0.90) -- (1.0,0.90) -- cycle;
  \draw[layergrey,thick] (0,0) -- (6.1,0) -- (7.1,0.90) -- (1.0,0.90) -- cycle;
  \fill[layerteal!12] (0.91,0.82) rectangle (1.43,1.38);
  \draw[layerteal!70,line width=0.4pt] (0.91,0.82) rectangle (1.43,1.38);
  \fill[layerteal!30] (0.91,1.38) -- (1.08,1.47) -- (1.60,1.47) -- (1.43,1.38) -- cycle;
  \draw[layerteal!70,line width=0.4pt] (0.91,1.38) -- (1.08,1.47) -- (1.60,1.47);
  \fill[layerteal!12] (1.67,0.82) rectangle (2.19,1.60);
  \draw[layerteal!70,line width=0.4pt] (1.67,0.82) rectangle (2.19,1.60);
  \fill[layerteal!30] (1.67,1.60) -- (1.83,1.69) -- (2.35,1.69) -- (2.19,1.60) -- cycle;
  \draw[layerteal!70,line width=0.4pt] (1.67,1.60) -- (1.83,1.69) -- (2.35,1.69);
  \fill[layerteal!12] (2.42,0.82) rectangle (2.93,1.97);
  \draw[layerteal!70,line width=0.4pt] (2.42,0.82) rectangle (2.93,1.97);
  \fill[layerteal!30] (2.42,1.97) -- (2.58,2.06) -- (3.10,2.06) -- (2.93,1.97) -- cycle;
  \draw[layerteal!70,line width=0.4pt] (2.42,1.97) -- (2.58,2.06) -- (3.10,2.06);
  \fill[layerteal!12] (3.17,0.82) rectangle (3.68,2.30);
  \draw[layerteal!70,line width=0.4pt] (3.17,0.82) rectangle (3.68,2.30);
  \fill[layerteal!30] (3.17,2.30) -- (3.33,2.39) -- (3.85,2.39) -- (3.68,2.30) -- cycle;
  \draw[layerteal!70,line width=0.4pt] (3.17,2.30) -- (3.33,2.39) -- (3.85,2.39);
  \fill[layerteal!12] (3.92,0.82) rectangle (4.43,2.32);
  \draw[layerteal!70,line width=0.4pt] (3.92,0.82) rectangle (4.43,2.32);
  \fill[layerteal!30] (3.92,2.32) -- (4.08,2.41) -- (4.60,2.41) -- (4.43,2.32) -- cycle;
  \draw[layerteal!70,line width=0.4pt] (3.92,2.32) -- (4.08,2.41) -- (4.60,2.41);
  \fill[layerteal!12] (4.67,0.82) rectangle (5.18,2.00);
  \draw[layerteal!70,line width=0.4pt] (4.67,0.82) rectangle (5.18,2.00);
  \fill[layerteal!30] (4.67,2.00) -- (4.83,2.09) -- (5.35,2.09) -- (5.18,2.00) -- cycle;
  \draw[layerteal!70,line width=0.4pt] (4.67,2.00) -- (4.83,2.09) -- (5.35,2.09);
  \fill[layerteal!12] (5.42,0.82) rectangle (5.93,1.62);
  \draw[layerteal!70,line width=0.4pt] (5.42,0.82) rectangle (5.93,1.62);
  \fill[layerteal!30] (5.42,1.62) -- (5.58,1.71) -- (6.10,1.71) -- (5.93,1.62) -- cycle;
  \draw[layerteal!70,line width=0.4pt] (5.42,1.62) -- (5.58,1.71) -- (6.10,1.71);
  \fill[layerteal!12] (6.17,0.82) rectangle (6.68,1.39);
  \draw[layerteal!70,line width=0.4pt] (6.17,0.82) rectangle (6.68,1.39);
  \fill[layerteal!30] (6.17,1.39) -- (6.33,1.48) -- (6.85,1.48) -- (6.68,1.39) -- cycle;
  \draw[layerteal!70,line width=0.4pt] (6.17,1.39) -- (6.33,1.48) -- (6.85,1.48);
  \fill[layerteal!12] (0.59,0.50) rectangle (1.11,1.07);
  \draw[layerteal!70,line width=0.4pt] (0.59,0.50) rectangle (1.11,1.07);
  \fill[layerteal!30] (0.59,1.07) -- (0.76,1.16) -- (1.28,1.16) -- (1.11,1.07) -- cycle;
  \draw[layerteal!70,line width=0.4pt] (0.59,1.07) -- (0.76,1.16) -- (1.28,1.16);
  \fill[layerteal!12] (1.34,0.50) rectangle (1.86,1.33);
  \draw[layerteal!70,line width=0.4pt] (1.34,0.50) rectangle (1.86,1.33);
  \fill[layerteal!30] (1.34,1.33) -- (1.51,1.42) -- (2.03,1.42) -- (1.86,1.33) -- cycle;
  \draw[layerteal!70,line width=0.4pt] (1.34,1.33) -- (1.51,1.42) -- (2.03,1.42);
  \fill[layerteal!12] (2.09,0.50) rectangle (2.61,1.75);
  \draw[layerteal!70,line width=0.4pt] (2.09,0.50) rectangle (2.61,1.75);
  \fill[layerteal!30] (2.09,1.75) -- (2.26,1.84) -- (2.78,1.84) -- (2.61,1.75) -- cycle;
  \draw[layerteal!70,line width=0.4pt] (2.09,1.75) -- (2.26,1.84) -- (2.78,1.84);
  \fill[layerteal!12] (2.84,0.50) rectangle (3.36,2.14);
  \draw[layerteal!70,line width=0.4pt] (2.84,0.50) rectangle (3.36,2.14);
  \fill[layerteal!30] (2.84,2.14) -- (3.01,2.23) -- (3.53,2.23) -- (3.36,2.14) -- cycle;
  \draw[layerteal!70,line width=0.4pt] (2.84,2.14) -- (3.01,2.23) -- (3.53,2.23);
  \fill[layerteal!12] (3.59,0.50) rectangle (4.11,2.15);
  \draw[layerteal!70,line width=0.4pt] (3.59,0.50) rectangle (4.11,2.15);
  \fill[layerteal!30] (3.59,2.15) -- (3.76,2.24) -- (4.28,2.24) -- (4.11,2.15) -- cycle;
  \draw[layerteal!70,line width=0.4pt] (3.59,2.15) -- (3.76,2.24) -- (4.28,2.24);
  \fill[layerteal!12] (4.34,0.50) rectangle (4.86,1.79);
  \draw[layerteal!70,line width=0.4pt] (4.34,0.50) rectangle (4.86,1.79);
  \fill[layerteal!30] (4.34,1.79) -- (4.51,1.88) -- (5.03,1.88) -- (4.86,1.79) -- cycle;
  \draw[layerteal!70,line width=0.4pt] (4.34,1.79) -- (4.51,1.88) -- (5.03,1.88);
  \fill[layerteal!12] (5.09,0.50) rectangle (5.61,1.35);
  \draw[layerteal!70,line width=0.4pt] (5.09,0.50) rectangle (5.61,1.35);
  \fill[layerteal!30] (5.09,1.35) -- (5.26,1.44) -- (5.78,1.44) -- (5.61,1.35) -- cycle;
  \draw[layerteal!70,line width=0.4pt] (5.09,1.35) -- (5.26,1.44) -- (5.78,1.44);
  \fill[layerteal!12] (5.84,0.50) rectangle (6.36,1.08);
  \draw[layerteal!70,line width=0.4pt] (5.84,0.50) rectangle (6.36,1.08);
  \fill[layerteal!30] (5.84,1.08) -- (6.01,1.17) -- (6.53,1.17) -- (6.36,1.08) -- cycle;
  \draw[layerteal!70,line width=0.4pt] (5.84,1.08) -- (6.01,1.17) -- (6.53,1.17);
  \fill[layerteal!12] (0.26,0.17) rectangle (0.78,0.73);
  \draw[layerteal!70,line width=0.4pt] (0.26,0.17) rectangle (0.78,0.73);
  \fill[layerteal!30] (0.26,0.73) -- (0.43,0.82) -- (0.95,0.82) -- (0.78,0.73) -- cycle;
  \draw[layerteal!70,line width=0.4pt] (0.26,0.73) -- (0.43,0.82) -- (0.95,0.82);
  \fill[layerteal!12] (1.02,0.17) rectangle (1.54,0.95);
  \draw[layerteal!70,line width=0.4pt] (1.02,0.17) rectangle (1.54,0.95);
  \fill[layerteal!30] (1.02,0.95) -- (1.19,1.04) -- (1.71,1.04) -- (1.54,0.95) -- cycle;
  \draw[layerteal!70,line width=0.4pt] (1.02,0.95) -- (1.19,1.04) -- (1.71,1.04);
  \fill[layerteal!12] (1.76,0.17) rectangle (2.29,1.32);
  \draw[layerteal!70,line width=0.4pt] (1.76,0.17) rectangle (2.29,1.32);
  \fill[layerteal!30] (1.76,1.32) -- (1.93,1.41) -- (2.46,1.41) -- (2.29,1.32) -- cycle;
  \draw[layerteal!70,line width=0.4pt] (1.76,1.32) -- (1.93,1.41) -- (2.46,1.41);
  \fill[layerteal!12] (2.51,0.17) rectangle (3.04,1.65);
  \draw[layerteal!70,line width=0.4pt] (2.51,0.17) rectangle (3.04,1.65);
  \fill[layerteal!30] (2.51,1.65) -- (2.69,1.74) -- (3.21,1.74) -- (3.04,1.65) -- cycle;
  \draw[layerteal!70,line width=0.4pt] (2.51,1.65) -- (2.69,1.74) -- (3.21,1.74);
  \fill[layerteal!12] (3.26,0.17) rectangle (3.79,1.67);
  \draw[layerteal!70,line width=0.4pt] (3.26,0.17) rectangle (3.79,1.67);
  \fill[layerteal!30] (3.26,1.67) -- (3.44,1.76) -- (3.96,1.76) -- (3.79,1.67) -- cycle;
  \draw[layerteal!70,line width=0.4pt] (3.26,1.67) -- (3.44,1.76) -- (3.96,1.76);
  \fill[layerteal!12] (4.01,0.17) rectangle (4.53,1.35);
  \draw[layerteal!70,line width=0.4pt] (4.01,0.17) rectangle (4.53,1.35);
  \fill[layerteal!30] (4.01,1.35) -- (4.18,1.44) -- (4.70,1.44) -- (4.53,1.35) -- cycle;
  \draw[layerteal!70,line width=0.4pt] (4.01,1.35) -- (4.18,1.44) -- (4.70,1.44);
  \fill[layerteal!12] (4.76,0.17) rectangle (5.28,0.97);
  \draw[layerteal!70,line width=0.4pt] (4.76,0.17) rectangle (5.28,0.97);
  \fill[layerteal!30] (4.76,0.97) -- (4.93,1.06) -- (5.45,1.06) -- (5.28,0.97) -- cycle;
  \draw[layerteal!70,line width=0.4pt] (4.76,0.97) -- (4.93,1.06) -- (5.45,1.06);
  \fill[layerteal!12] (5.51,0.17) rectangle (6.03,0.74);
  \draw[layerteal!70,line width=0.4pt] (5.51,0.17) rectangle (6.03,0.74);
  \fill[layerteal!30] (5.51,0.74) -- (5.68,0.83) -- (6.20,0.83) -- (6.03,0.74) -- cycle;
  \draw[layerteal!70,line width=0.4pt] (5.51,0.74) -- (5.68,0.83) -- (6.20,0.83);
  \fill[accent!42] (2.69,0.17) rectangle (2.86,1.65);
  \draw[accent,line width=1pt] (2.69,0.17) rectangle (2.86,1.65);
  \fill[accent!55] (2.69,1.65) -- (2.78,1.70) -- (2.96,1.70) -- (2.86,1.65) -- cycle;
  \draw[accent,line width=1pt] (2.69,1.65) -- (2.78,1.70) -- (2.96,1.70);
  \draw[accent,densely dashed] (2.86,0.16) -- (3.42,-0.26);
  \node[accent,anchor=north west,align=left]
    at (3.50,-0.30) {one thin column: base $\Delta A$,\\ height $f(x,y)$, volume $\approx f(x,y)\,\Delta A$};
  \node[layergrey,anchor=north west] at (-0.05,-0.08) {base $R$};
  \node[layergrey,anchor=north] at (5.95,-0.02) {$x$};
  \node[layergrey,anchor=west] at (7.15,0.55) {$y$};
\end{tikzpicture}
\caption{The solid under the graph of $f$ over a rectangle $R$, cut into
thin columns. Each column is almost a box with base $\Delta A$ and height
$f(x,y)$, so its volume is $f(x,y)\,\Delta A$; adding up all the columns gives
$\iint_R f\,dA$. The column marked in red is one term in that sum.
Slicing along the $x$ direction and along the $y$ direction gives different columns, but
the total volume is the same.}
\label{fig:kolom-tipis}
\end{figure}

The regions that appear in problems are rarely rectangles. We handle them
with one extra step: extend the function by zero outside the region.

\begin{lemma}[A region between two graphs]\index{type I region}
Let $g_{1},g_{2}\colon[a,b]\to\mathbb{R}$ be continuous with $g_{1}\le
g_{2}$, and
\[
D=\{(x,y): a\le x\le b,\ g_{1}(x)\le y\le g_{2}(x)\}.
\]
If $f$ is continuous on $D$, then
\[
\iint_{D}f\,dA=\int_{a}^{b}\!\!\int_{g_{1}(x)}^{g_{2}(x)}f(x,y)\,dy\,dx .
\]
\end{lemma}

\begin{proof}
Take a rectangle $R$ that contains $D$, and extend $f$ by the value $0$ on
$R\setminus D$. The boundary of $D$ can be covered by finitely many rectangles of
total area less than $\varepsilon$: on the graph, uniform continuity on a
compact interval gives a partition whose oscillation of $g$ is small on each piece.
A set that can be covered that way does not change the value of the integral, so
\[
\iint_{D}f\,dA=\iint_{R} f\cdot\mathbf{1}_{D}\,dA .
\]
By Fubini, the integral on the right is computed by integrating $x$ first;
for fixed $x$, the integrand $\mathbf{1}_{D}f$ equals $f$ only on
$g_{1}(x)\le y\le g_{2}(x)$.
\end{proof}

A type I region is enough for almost every problem in this chapter; a boundary that crosses a
vertical line more than twice is split and then added up. The most important case is the disk:
in Cartesian coordinates it demands two pieces, in polar only one.

\begin{untukanak}
Now take a closed room with several stoves inside it.
Each stove gives off heat. The question: how much heat leaves that room
every second, added up over all the walls, the floor, and the roof?

The first way to compute it is hard: walk around every wall, measure the heat
flow at each point, then add up. The second way is much easier: count
how much heat all the stoves inside produce. If no heat escapes
through other gaps, the two numbers must agree: the amount inside equals the
amount through the skin.

That sentence sounds too good to be true. The agreement holds only if
the skin of the room is really closed and the direction ``out'' is fixed unambiguously. If one small
hole is not counted, or if we mislabel which side is outside and
which is inside, the numbers miss even though all our additions are correct.

In this chapter we shall learn why the amount inside can be read off from the skin,
for any curved shape, and how one small hole can ruin
the whole conclusion.
\end{untukanak}

\section{Change of variables and the Jacobian determinant}
\tingkat{intermediate}

Changing coordinates means changing cells: a square cell in the $(u,v)$ plane
becomes a curved cell in the $(x,y)$ plane, and its area changes. We need
one number at each point.

\begin{definisi}[Change of variables]
\index{change of variables}\index{Jacobian} Let $D^{*}$ be an open region in
$\mathbb{R}^{2}$ and $T\colon D^{*}\to\mathbb{R}^{2}$,
$T(u,v)=(x(u,v),y(u,v))$, a function with continuous partial derivatives. $T$ is called a
\emph{valid change of variables} if $T$ is one-to-one on the interior of
$D^{*}$ and its Jacobian determinant
\[
\det J_{T}(u,v)=\frac{\partial x}{\partial u}\frac{\partial y}{\partial v}
-\frac{\partial x}{\partial v}\frac{\partial y}{\partial u}
\]
is nonzero there. The matrix $J_{T}$ itself we already met in Chapter 12.
\end{definisi}

In everyday words: it is the number that tells how many times a small cell
grows when it is carried by $T$. A value of $3$ means the cell becomes three times larger in area, $1/2$ means it shrinks
by half, and zero means the cell is flattened into a line and its information is lost.

Figure~\ref{fig:juring-jacobian} shows one cell and its image:
area $0.04$ becomes $0.064$, and the number $1.6$ is $|{\det J}|$ at the center
of the cell.

\begin{figure}[htbp]
\centering
\begin{tikzpicture}[font=\footnotesize]
  \begin{scope}[x=1.9cm,y=1.9cm]
    \fill[layergrey!7] (1,0) rectangle (2.4,1.2);
    \draw[layergrey!40,very thin] (1,0) grid[xstep=0.2,ystep=0.2] (2.4,1.2);
    \draw[->,ink] (0.88,0) -- (2.62,0) node[right] {$u$};
    \draw[->,ink] (1,-0.08) -- (1,1.42) node[above] {$v$};
    \fill[layerindigo!30] (1.5,0.6) rectangle (1.7,0.8);
    \draw[layerindigo,thick] (1.5,0.6) rectangle (1.7,0.8);
    \node[layerindigo!85!black,anchor=north,align=center,fill=white,inner sep=1.2pt]
      at (1.60,0.55) {cell $0.2\times0.2$\\ area $0.04$};
    \node[layergrey,anchor=north] at (1.7,-0.10) {the $(u,v)$ plane};
  \end{scope}
  \draw[-{Latex[length=2.4mm]},ink,thick] (4.62,1.14) -- (5.85,1.14);
  \node[ink,above] at (5.23,1.18) {$T$};
  \node[ink,below,font=\scriptsize] at (5.23,1.06) {$\times\,1.6$};
  \begin{scope}[xshift=5.9cm,x=1.9cm,y=1.9cm]
    \foreach \r in {1,1.2,1.4,1.6,1.8,2,2.2,2.4}{
      \draw[layergrey!40,very thin] (\r,0) arc (0:69:\r);}
    \foreach \a in {0,11.5,23,34.5,46,57.5,69}{
      \draw[layergrey!40,very thin] (0,0) -- (\a:2.4);}
    \draw[->,ink] (-0.08,0) -- (2.62,0) node[right] {$x$};
    \draw[->,ink] (0,-0.08) -- (0,2.52) node[above] {$y$};
    \fill[accent!30] (34.4:1.5) arc (34.4:45.8:1.5) -- (45.8:1.7)
      arc (45.8:34.4:1.7) -- cycle;
    \draw[accent,thick] (34.4:1.5) arc (34.4:45.8:1.5) -- (45.8:1.7)
      arc (45.8:34.4:1.7) -- cycle;
    \node[accent,anchor=west,align=left] at (1.62,1.42)
      {curved sector\\ area $0.064=1.6\times0.04$};
    \node[layergrey,anchor=north] at (1.6,-0.10) {the $(x,y)$ plane};
  \end{scope}
\end{tikzpicture}
\caption{One small cell in the $(u,v)$ plane and its image in the $(x,y)$ plane under
$T(u,v)=(u\cos v,\;u\sin v)$. The cell $0.2\times0.2$ of area $0.04$ becomes a sector of
area $0.064$, that is $1.6$ times as large, exactly the value $|{\det J}|=u$ at the center of the cell
($u=1.6$). Note that its sides curve as well: what changes is not only the size
but also the shape, and one number at each point is enough to record the change.}
\label{fig:juring-jacobian}
\end{figure}

\begin{tangan}
Take a square sheet of paper $12$ cm by $12$ cm, a ruler, a pencil, scissors, and
tape. Draw $1$ cm cells over the whole sheet; there are $144$ cells. Cut the paper
into twelve vertical strips, each $1$ cm by $12$ cm. On every
strip, mark the point $3$ cm to the right on the top edge, draw a line from the bottom left
corner to that point, then cut along the line. Each strip is now two slanted
pieces. Arrange all the pieces into one tight parallelogram, then
tape them together.

Watch what happens: the outer frame is $15$ cm by $12$ cm, wider by $3$ cm
than the original sheet, but the parallelogram inside still has base
$12$ cm and height $12$ cm. Measure both with the ruler, perpendicular, not
along the slanted side: the result is $144$ square centimeters, the same as the area of the sheet
before it was cut. If your number is different, the height is almost certainly what went wrong.
\end{tangan}

\begin{teorema}[Change of variables in a double integral]
\index{change of variables theorem} Let $T$ be a valid change of variables
from $D^{*}$ to $D=T(D^{*})$, and $f$ continuous on $D$. Then
\[
\iint_{D}f(x,y)\,dx\,dy
=\iint_{D^{*}} f\bigl(T(u,v)\bigr)\,\bigl|\det J_{T}(u,v)\bigr|\,du\,dv .
\]
\end{teorema}

\begin{proof}
We prove the idea; the full proof is in Spivak, the chapter \emph{Integration},
the section on change of variables, and in Rudin, the chapter \emph{Integration of Differential
Forms}.

Partition $D^{*}$ into small cells. On a cell around $(u_{0},v_{0})$,
the map $T$ is almost linear: $T(u_{0}+h,v_{0}+k)\approx
T(u_{0},v_{0})+J_{T}(u_{0},v_{0})(h,k)$. That cell is mapped onto a
parallelogram with sides $h\,\partial T/\partial u$ and $k\,\partial T/\partial v$,
whose area is computed with a determinant (Chapter 7):
\[
\mathrm{luas}
=\Bigl|h\,k\,\det\Bigl(\frac{\partial T}{\partial u},
\frac{\partial T}{\partial v}\Bigr)\Bigr|
=|h\,k|\,|\det J_{T}(u_{0},v_{0})| .
\]
Each cell contributes $f(T(u_{0},v_{0}))$ times its original area times
$|\det J_{T}|$, and the sum is the Riemann sum for the integral on the right-hand
side; because the derivative of $T$ is continuous on a compact set, the error of the linear
approximation shrinks uniformly with the mesh.
\end{proof}

The absolute value on the determinant is not decoration: a determinant can be negative, an area cannot.
The negative sign stores an orientation; in the double integral that orientation is discarded, but
in the line and surface integrals later that sign decides the answer.

Now the determinants for the two coordinate systems used most often,
beginning with polar.

\begin{proposisi}[Polar coordinates]\index{polar coordinates}
With $x=r\cos\theta$, $y=r\sin\theta$, $r>0$, we have
\[
\det J=\begin{vmatrix}\cos\theta & -r\sin\theta\\ \sin\theta & r\cos\theta
\end{vmatrix}=r\cos^{2}\theta+r\sin^{2}\theta=r ,
\]
so that
$\iint_{D}f\,dx\,dy=\iint_{D^{*}}f(r\cos\theta,r\sin\theta)\,r\,dr\,d\theta$.
\end{proposisi}

That factor $r$ is the one missing from the program in the opening of the chapter. At $r=1$
it is $1$, but at $r=0.1$ it is $0.1$: a cell near the center
contributes ten times less than a guess that ignores it, and
the total error can be more than a factor of two.

\begin{proposisi}[Cylindrical and spherical coordinates]
\index{cylindrical coordinates}\index{spherical coordinates} For cylindrical coordinates
$x=r\cos\theta$, $y=r\sin\theta$, $z=z$, the Jacobian determinant is $r$. For
spherical coordinates
\[
x=\rho\sin\varphi\cos\theta,\quad
y=\rho\sin\varphi\sin\theta,\quad
z=\rho\cos\varphi ,
\]
with $\rho\ge0$, $0\le\varphi\le\pi$, $0\le\theta<2\pi$, the Jacobian
determinant is $\rho^{2}\sin\varphi$.
\end{proposisi}

\begin{proof}
For the cylinder, the third column $(0,0,1)$ leaves the determinant as the $2\times2$
polar block determinant, namely $r$.

For the sphere, compute the $3\times3$ determinant by cofactor expansion. With
$s_{\varphi}=\sin\varphi$, $c_{\varphi}=\cos\varphi$,
$s_{\theta}=\sin\theta$, $c_{\theta}=\cos\theta$, the matrix is
\[
J=\begin{pmatrix}
s_{\varphi}c_{\theta} & \rho c_{\varphi}c_{\theta} & -\rho s_{\varphi}s_{\theta}\\
s_{\varphi}s_{\theta} & \rho c_{\varphi}s_{\theta} & \rho s_{\varphi}c_{\theta}\\
c_{\varphi} & -\rho s_{\varphi} & 0
\end{pmatrix}.
\]
The entry in the third row and third column is zero, so expansion along the third column leaves two
terms:
\[
\det J=\rho^{2}s_{\varphi}s_{\theta}^{2}+\rho^{2}s_{\varphi}c_{\theta}^{2}
=\rho^{2}s_{\varphi}\bigl(s_{\theta}^{2}+c_{\theta}^{2}\bigr)=\rho^{2}s_{\varphi}.
\]
\end{proof}

The factor $\rho^{2}\sin\varphi$ is zero at the poles $\varphi=0$ and $\varphi=\pi$,
because there the whole sphere of radius $\rho$ collapses to a single point and the map stops being
one-to-one; that is why the poles are excluded from $D^{*}$. It is largest at the
equator, where a small change in $\varphi$ carries a point farthest.

\begin{resep}
\begin{enumerate}
\item Write $T$ explicitly: $x=x(u,v)$, $y=y(u,v)$, and check that $T$
  is one-to-one on the interior of $D^{*}$.
\item Compute the four partial derivatives, assemble the $2\times2$ matrix, then compute
  $\det J(u,v)$. Keep its sign; the absolute value is put in later.
\item Translate the boundary: draw $D$ in the $(x,y)$ plane, ask which values of $(u,v)$
  land there, then write $D^{*}$.
\item Replace every $x,y$ in $f$ by $x(u,v),y(u,v)$, multiply by
  $|\det J(u,v)|$, then close with $du\,dv$ in the order of the limits.
\end{enumerate}
A warning about order: steps 2 and 3 may not be swapped. The determinant must already
be in terms of $(u,v)$ before the limits are put in. If the limits are written first, that factor is left in $x$ and $y$ and
the integral cannot be computed.
\end{resep}

\kamus{Jacobian determinant}{determinan Jacobian}{The number that says how many times a small cell stretches when the coordinates are changed.}

\kamus{double integral}{integral lipat dua}{Adding up values on small cells of a surface, then narrowing the cells.}

\kamus{orientation}{orientasi}{An agreement about which direction of rotation counts as positive.}

\section{Line integrals and conservative fields}
\tingkat{intermediate}

Now we stop adding up over a surface and start adding up
along a path. Picture a small object dragged along a curved path by a force whose
direction and size change at every point. The work of that force is the sum of the component of the force along the path times the piece of
length of the path.

\begin{definisi}[Line integral]\index{line integral}
Let $\gamma\colon[a,b]\to\mathbb{R}^{2}$ be a curve with piecewise continuous
derivative, and $F=(P,Q)$ a continuous vector field on its trace.
The \emph{line integral} of $F$ along $\gamma$ is
\[
\int_{\gamma}F\cdot d\mathbf{r}
=\int_{a}^{b}\Bigl(P(\gamma(t))\,\gamma_{1}'(t)
+Q(\gamma(t))\,\gamma_{2}'(t)\Bigr)\,dt .
\]
If $\gamma$ is closed, meaning $\gamma(a)=\gamma(b)$, that integral is called
the \emph{circulation} of $F$ around $\gamma$.
\end{definisi}

That formula does not choose a starting point, only a direction. Changing the parameter to a
different speed does not change the value, because substitution (Chapter 13) absorbs
the speed factor; reversing the direction changes the sign. So a line integral is oriented,
and that is its first difference from the double integral.

Figure~\ref{fig:sirkulasi} shows such a rotating field, with one
closed curve for measuring how much rotation is caught.

\begin{figure}[htbp]
\centering
\begin{tikzpicture}[scale=1.45,font=\footnotesize]
  \draw[->,layergrey] (-2.85,0) -- (2.85,0) node[right] {$x$};
  \draw[->,layergrey] (0,-2.85) -- (0,2.85) node[above] {$y$};
  \foreach \x in {-2.4,-1.6,-0.8,0,0.8,1.6,2.4}{
    \foreach \y in {-2.4,-1.6,-0.8,0,0.8,1.6,2.4}{
      \draw[-{Latex[length=1.2mm]},layerteal!75]
        (\x,\y) -- ({\x-0.30*\y},{\y+0.30*\x});}}
  \draw[accent,very thick] (0,0) circle (1.75);
  \draw[-{Latex[length=2.8mm]},accent,very thick]
    (1.126,-1.341) -- (1.340,-1.126);
  \node[accent,anchor=north west,align=left] at (0.28,-1.86)
    {closed curve, positive direction};
  \node[ink,align=center,fill=white,fill opacity=0.85,text opacity=1,inner sep=2pt] at (0,0.95)
    {$\displaystyle\oint_{\gamma}F\cdot d\mathbf{r}=2\pi$};
  \node[ink,font=\scriptsize,fill=white,fill opacity=0.85,text opacity=1,inner sep=1.5pt] at (0,0.42)
    {whereas $\operatorname{curl}F=2$ at every point};
\end{tikzpicture}
\caption{The field $F(x,y)=(-y,x)$ with an arrow at every grid point. The arrows turn
around the origin, and their length grows with the distance of the point from the center. The circulation
around a circle of radius $1$ is $2\pi$, equal to the area $\pi$ times
$\operatorname{curl}F=2$: this is Green for this field. The direction of the arrows on the curve is the positive
direction, counterclockwise.}
\label{fig:sirkulasi}
\end{figure}

If a force field comes from a height function, the work depends only
on the starting point and the ending point. This is the curve version of the fundamental theorem of calculus,
and its name follows from that.

\begin{teorema}[The fundamental theorem for line integrals]
\index{conservative field} Let $f\colon U\to\mathbb{R}$ have continuous partial
derivatives on an open set $U\subseteq\mathbb{R}^{2}$, and let $\gamma$
be a curve from $A$ to $B$ inside $U$. Then
\[
\int_{\gamma}\nabla f\cdot d\mathbf{r}=f(B)-f(A).
\]
\end{teorema}

\begin{proof}
Write $g(t)=f(\gamma(t))$. By the multivariable chain rule (Chapter 12),
\[
g'(t)=\nabla f(\gamma(t))\cdot\gamma'(t),
\]
provided $\gamma$ is differentiable at $t$. So the integrand in the line integral
is the derivative of $g$, and the fundamental theorem of calculus (Chapter 13) gives
$\int_{a}^{b}g'(t)\,dt=g(b)-g(a)=f(B)-f(A)$. For a curve that is only
piecewise differentiable, add up over each piece; the terms at the joining
points cancel each other.
\end{proof}

A field of the form $\nabla f$ is called \emph{conservative}, and $f$ is its potential.
A direct consequence of the theorem: for a conservative field, the circulation around
a closed path is always zero.

The converse is also true.

\begin{teorema}[Converse]
Let $F$ be continuous on a connected open set
$U\subseteq\mathbb{R}^{2}$, and let the line integral of $F$ be independent of the
path inside $U$. Then $F=\nabla f$ for some $f$, namely
\[
f(\mathbf{x})=\int_{\mathbf{x}_{0}}^{\mathbf{x}}F\cdot d\mathbf{r},
\]
with $\mathbf{x}_{0}\in U$ fixed.
\end{teorema}

\begin{proof}
Take $\mathbf{x}\in U$ and $h$ small enough that the segment from $\mathbf{x}$ to
$\mathbf{x}+h\mathbf{e}_{1}$ stays inside $U$. Because the path does not
matter, we may travel from $\mathbf{x}_{0}$ to $\mathbf{x}$ and then
add that short segment, so that
\[
f(\mathbf{x}+h\mathbf{e}_{1})-f(\mathbf{x})
=\int_{0}^{1}F_{1}(\mathbf{x}+th\mathbf{e}_{1})\,h\,dt .
\]
Divide by $h$ and use continuity: the average value of the integrand tends to
$F_{1}(\mathbf{x})$ as $h\to0$, so $\partial f/\partial x_{1}=F_{1}$,
and the same argument gives $\partial f/\partial x_{2}=F_{2}$. Both partial
derivatives are continuous, hence $F=\nabla f$.
\end{proof}

Given $F$, how do we know it is conservative without looking for a potential? In
the plane there is a quick test: $F=(P,Q)$ has a potential only if $\partial P/\partial
y=\partial Q/\partial x$, because both sides are second derivatives of $f$ and Clairaut
(Chapter 12) makes them equal. The other direction demands a domain without holes and is
proved through Green later. The field from the opening of the chapter is the counterexample: its curl is zero,
but it is not the gradient of anything.

\begin{contoh}[Two ways to compute work]
\index{potential} Take $F(x,y)=(2xy,\ x^{2}+3y^{2})$ and the path
$\gamma(t)=(t,2t^{2})$ for $0\le t\le1$, from $(0,0)$ to $(1,2)$.

First way, directly. Here $dx=dt$ and $dy=4t\,dt$, so
\[
\int_{\gamma}F\cdot d\mathbf{r}
=\int_{0}^{1}\Bigl(2t\cdot 2t^{2}+(t^{2}+12t^{4})\cdot 4t\Bigr)dt
=\int_{0}^{1}\bigl(8t^{3}+48t^{5}\bigr)dt=2+8=10 .
\]
Second way, through the potential. First check that $\partial P/\partial
y=2x=\partial Q/\partial x$, then guess $f(x,y)=x^{2}y+y^{3}$. Its derivatives
are indeed $\partial f/\partial x=2xy$ and $\partial f/\partial y=x^{2}+3y^{2}$.
Therefore
\[
\int_{\gamma}F\cdot d\mathbf{r}=f(1,2)-f(0,0)=(2+8)-0=10 .
\]
Both ways give $10$; adaptive quadrature differs by $1.78\times10^{-15}$.
\end{contoh}

\section{Green's theorem}
\tingkat{intermediate}

The idea in the box about the room and the stoves now becomes a theorem: what happens
inside a region can be read off from its edge.

Before that, one agreement on direction.

\begin{definisi}[Positive orientation]
\index{positive orientation} A simple closed curve $\gamma$ is oriented
\emph{positively} if the region it bounds always lies on the left
as we traverse it. For a circle, the positive direction is counterclockwise.
\end{definisi}

\begin{teorema}[Green]\index{Green's theorem}
Let $D$ be a closed bounded region in $\mathbb{R}^{2}$ whose boundary
$\partial D$ consists of finitely many simple closed curves that are piecewise
smooth and positively oriented. Let $P,Q$ have continuous partial
derivatives on an open set containing $D$. Then
\[
\oint_{\partial D}P\,dx+Q\,dy=\iint_{D}\Bigl(\frac{\partial Q}{\partial x}
-\frac{\partial P}{\partial y}\Bigr)dA .
\]
\end{teorema}

\begin{proof}
It is enough to prove it for regions that are both type I and type II, because a general
region is split into finitely many pieces of that kind and the integrals are
added up; the adjacent sides of the pieces are traversed in opposite directions and therefore
cancel each other.

Suppose $D$ is type I with $g_{1},g_{2}$ smooth. On the vertical sides $dx=0$,
so their contribution is zero and
\[
\oint_{\partial D}P\,dx=\int_{a}^{b}P(x,g_{1}(x))\,dx+\int_{b}^{a}P(x,g_{2}(x))\,dx
=\int_{a}^{b}\bigl[P(x,g_{1})-P(x,g_{2})\bigr]dx .
\]
By the fundamental theorem of calculus,
$P(x,g_{1})-P(x,g_{2})=-\int_{g_{1}(x)}^{g_{2}(x)}\partial P/\partial y\,dy$,
so that
\[
\oint_{\partial D}P\,dx=-\int_{a}^{b}\!\!\int_{g_{1}(x)}^{g_{2}(x)}
\frac{\partial P}{\partial y}\,dy\,dx=-\iint_{D}\frac{\partial P}{\partial y}\,dA .
\]
In the same way, on a $D$ of type II, where $dy=0$ on the flat sides,
$\oint_{\partial D}Q\,dy=\iint_{D}\partial Q/\partial x\,dA$.
\end{proof}

\begin{contoh}[Green on the unit square]
Take $P=x^{2}-y$, $Q=x+2y$, and $D=[0,1]\times[0,1]$. Here $\partial
Q/\partial x-\partial P/\partial y=1-(-1)=2$, so the right-hand side of Green is
$2$ times the area of the square, namely $2$.

We compute the left-hand side on the four sides, counterclockwise:
\begin{align*}
\text{bottom }(y=0,\,dy=0):&\quad \int_{0}^{1}x^{2}\,dx=\tfrac13 ,\\
\text{right }(x=1,\,dx=0):&\quad \int_{0}^{1}(1+2y)\,dy=2 ,\\
\text{top }(y=1,\,dy=0):&\quad \int_{1}^{0}(x^{2}-1)\,dx=\tfrac23 ,\\
\text{left }(x=0,\,dx=0):&\quad \int_{1}^{0}2y\,dy=-1 .
\end{align*}
The sum is $\tfrac13+2+\tfrac23-1=2$, the same as the right-hand side. Adaptive quadrature
for the four sides gives $2.000000000000000$, and the double integral also gives
$2.000000000000000$.
\end{contoh}

Now the sign error we asked for. Repeat the same computation clockwise. Every side is traversed
backwards, so every term changes sign and
the sum becomes $-2$; the right-hand side of Green does not change because the region is the same.
The answer that should be $2$ becomes $-2$: not slightly off, but wrong in
sign entirely, with a difference of $4$ that cannot be blamed on
rounding.

So orientation is not a polite convention: Green's formula is correct only for one
direction of rotation, and that direction must be checked before computing.

With Green in hand, the converse direction of the curl test from the previous section
can be proved.

\begin{akibat}[The curl test on a domain without holes]\index{curl test}
Let $U\subseteq\mathbb{R}^{2}$ be open and \emph{simply connected},
meaning that every closed curve in $U$ can be contracted to a single point inside
$U$. Let $P,Q$ have continuous partial derivatives on $U$ and $\partial
P/\partial y=\partial Q/\partial x$ throughout $U$. Then $F=(P,Q)$ is conservative
on $U$.
\end{akibat}

\begin{proof}
Take a simple closed curve $\gamma$ in $U$ and the region $D$ that
it bounds. Because $U$ is simply connected, $D$ stays inside $U$, so
Green applies and
\begin{align*}
\oint_{\gamma}P\,dx+Q\,dy
&=\iint_{D}\Bigl(\frac{\partial Q}{\partial x}-\frac{\partial P}{\partial y}\Bigr)dA\\
&=0 .
\end{align*}
A closed curve that crosses itself is split into simple arcs,
and each piece gives zero. All circulations are zero, so the line
integral is independent of path, and the converse theorem above supplies a potential.
\end{proof}

That hypothesis cannot be dropped. On a region with a hole, Green demands that the inner boundary
also be counted; that is where the field from the opening of the chapter escapes the curl test.

\section{Gauss's and Stokes' theorems}
\tingkat{advanced}

Green talks about a region in the plane and its edge. In three-dimensional space
there are two versions, and both use the divergence and the curl from Chapter 12.

\begin{definisi}[Divergence and curl]
\index{divergence}\index{curl} For a field $F=(F_{1},F_{2},F_{3})$ with
continuous partial derivatives,
\[
\operatorname{div}F=\frac{\partial F_{1}}{\partial x}
+\frac{\partial F_{2}}{\partial y}+\frac{\partial F_{3}}{\partial z},
\qquad
\operatorname{curl}F=\Bigl(\frac{\partial F_{3}}{\partial y}
-\frac{\partial F_{2}}{\partial z},\
\frac{\partial F_{1}}{\partial z}-\frac{\partial F_{3}}{\partial x},\
\frac{\partial F_{2}}{\partial x}-\frac{\partial F_{1}}{\partial y}\Bigr).
\]
\end{definisi}

The divergence measures how much field is born at a point: positive means a
source, negative a sink. The curl measures the rotation of the field around that point.

\begin{teorema}[Gauss, or the divergence theorem]\index{Gauss's theorem}
Let $E\subseteq\mathbb{R}^{3}$ be a closed bounded region whose boundary
$\partial E$ is a piecewise smooth surface, and let $\mathbf{n}$ be the unit
normal vector pointing outward. If $F$ has continuous partial derivatives
on an open set containing $E$, then
\[
\iiint_{E}\operatorname{div}F\,dV=\iint_{\partial E}F\cdot\mathbf{n}\,dA .
\]
\end{teorema}

\begin{proof}
We prove it for a box first.

Take $E=[a,b]\times[c,d]\times[e,f]$. For the third term,
\[
\iiint_{E}\frac{\partial F_{3}}{\partial z}\,dV
=\iint_{[a,b]\times[c,d]}\bigl[F_{3}(x,y,f)-F_{3}(x,y,e)\bigr]dx\,dy ,
\]
by Fubini and the fundamental theorem of calculus in the $z$ direction. That right-hand side is exactly
the flux through the top face (normal $+\mathbf{k}$) and the bottom face (normal
$-\mathbf{k}$), because there $\mathbf{n}\,dA=\pm\mathbf{k}\,dx\,dy$. The other
terms give the flux through the two other pairs of faces.

For a general region, cut it into small boxes. An interior face appears twice
with opposite normals and cancels, so all that remains are the faces
that touch the outer boundary; the limit as the mesh goes to zero gives the statement for a curved
boundary. The full proof is in Spivak, Theorem 5-8.
\end{proof}

\begin{contoh}[Gauss on the unit ball]
Take $F(x,y,z)=(x^{3},y^{3},z^{3})$ and $E$ the unit ball. Here
$\operatorname{div}F=3(x^{2}+y^{2}+z^{2})=3\rho^{2}$, so with the
Jacobian determinant of the sphere
\[
\iiint_{E}3\rho^{2}\,dV=\int_{0}^{2\pi}\!\!\int_{0}^{\pi}\!\!\int_{0}^{1}
3\rho^{4}\sin\varphi\,d\rho\,d\varphi\,d\theta
=3\cdot2\pi\cdot2\cdot\frac15=\frac{12\pi}{5}.
\]
The surface side: on the unit ball $\mathbf{n}=(x,y,z)$, so
$F\cdot\mathbf{n}=x^{4}+y^{4}+z^{4}$ and
\[
\iint_{\partial E}(x^{4}+y^{4}+z^{4})\,dA=3\iint_{\partial E}z^{4}\,dA
=3\int_{0}^{2\pi}\!\!\int_{0}^{\pi}\cos^{4}\varphi\sin\varphi\,d\varphi\,d\theta
=3\cdot2\pi\cdot\frac25=\frac{12\pi}{5},
\]
with the factor $3$ from symmetry. The value is $7.5398223686155$, and quadrature on
the surface gives the same number to the last digit.
\end{contoh}

The second theorem talks about an open surface whose boundary is a curve.
Its formula requires the right-hand rule.

\begin{teorema}[Stokes]\index{Stokes' theorem}
Let $S$ be an oriented piecewise smooth surface with normal
$\mathbf{n}$, and let $\partial S$ be its boundary curve with the orientation given by the
right-hand rule. If $F$ has continuous partial derivatives on an open set
containing $S$, then
\[
\iint_{S}(\operatorname{curl}F)\cdot\mathbf{n}\,dA
=\oint_{\partial S}F\cdot d\mathbf{r}.
\]
\end{teorema}

\begin{proof}
Parametrize $S$ by $\mathbf{r}(u,v)$ over a region $D$ in the $(u,v)$ plane,
with the matching orientation. The chain rule turns the integrand into
derivatives of the components of $F$ with respect to $u$ and $v$, and after rearrangement the surface
integral becomes
\[
\iint_{D}\Bigl(\frac{\partial \tilde{Q}}{\partial u}
-\frac{\partial \tilde{P}}{\partial v}\Bigr)du\,dv ,
\]
for two functions $\tilde{P},\tilde{Q}$ built from $F$ and
$\mathbf{r}$. That is exactly the right-hand side of Green, so Green gives the line integral
on $\partial D$, which is nothing other than the line integral of $F$ on $\partial S$.
The full algebra is in Spivak, Theorem 5-9.
\end{proof}

If $S$ is flat in the $xy$ plane, Stokes becomes Green; if $S$ is only an
interval, it becomes the fundamental theorem of calculus.

\begin{contoh}[Stokes on a hemisphere]
Take $F=(z,x,y)$, whose curl is $(1,1,1)$, and $S$ the upper half of the unit ball
with outward normal, so that its boundary is the unit circle in the $xy$ plane
counterclockwise.

The surface side:
\[
\iint_{S}(1,1,1)\cdot\mathbf{n}\,dA=\iint_{S}(x+y+z)\,dA=\iint_{S}z\,dA
=\int_{0}^{2\pi}\!\!\int_{0}^{\pi/2}\cos\varphi\sin\varphi\,d\varphi\,d\theta
=2\pi\cdot\frac12=\pi ,
\]
because the $x$ and $y$ terms cancel each other by the symmetry of the hemisphere. The boundary side, with
$\mathbf{r}(t)=(\cos t,\sin t,0)$:
\[
\oint_{\partial S}F\cdot d\mathbf{r}=\int_{0}^{2\pi}\cos^{2}t\,dt=\pi .
\]
Both are $3.1415926535898$; quadrature gives the same number on the two sides.
\end{contoh}

Now the sign error in Stokes. Turn the normal of the surface inward:
the surface side becomes $-\pi$, the boundary side does not change as long as the direction of rotation
stays the same, and the equality fails with a difference of $2\pi$. The right-hand rule is what
binds the two sides; reversing both together restores the equality.

\section{Differential forms: one sentence for all of them}
\tingkat{advanced}

We have now seen three theorems with three faces. There is a language that
writes them as one statement, and that language is not mere tidiness: it
explains why the sign of orientation appears everywhere.

\begin{definisi}[Differential forms]
\index{differential form} A \emph{$k$-form} on an open set
$U\subseteq\mathbb{R}^{n}$ is a sum of terms of the form
$a(x)\,dx_{i_{1}}\wedge\cdots\wedge dx_{i_{k}}$, with $a$ a smooth function.
The wedge product $\wedge$ is antisymmetric:
\[
dx\wedge dy=-dy\wedge dx,\qquad dx\wedge dx=0 .
\]
The exterior derivative $d$ acts on the coefficients and adds one wedge:
\[
d(a\,dx_{I})=da\wedge dx_{I},\qquad da=\sum_{j}\frac{\partial a}{\partial x_{j}}dx_{j}.
\]
\end{definisi}

In everyday words: a differential form is an object that is ready to be
integrated over a surface of a given dimension, and $d$ raises the dimension by one
level. Antisymmetry makes reversing the orientation reverse the sign automatically;
the sign difficulties from before are built into that one rule.

\begin{teorema}[General Stokes]
Let $M$ be an oriented manifold of dimension $k$ with boundary, and let $\omega$
be a $(k-1)$-form with smooth coefficients. Then
\[
\int_{\partial M}\omega=\int_{M}d\omega ,
\]
with the orientation of $\partial M$ induced by the orientation of $M$.
\end{teorema}

The complete proof is skipped: Spivak, \emph{Integration on Chains}, Theorem 4-13
for chains, and Theorem 5-5 in \emph{Integration of Manifolds} for
manifolds; also Rudin, the chapter \emph{Integration of Differential Forms}. There
Green, Gauss, and classical Stokes are derived as special cases.

Its translation in three-dimensional space:

\begin{center}
\begin{tabular}{lll}
\hline
Language of forms & Vector field & Theorem \\
\hline
0-form $f$ & scalar function & $df$ gives the gradient \\
1-form $P\,dx+Q\,dy+R\,dz$ & $F=(P,Q,R)$ & $d\omega$ gives $\operatorname{curl}F$ \\
2-form $P\,dy\wedge dz+\cdots$ & $F=(P,Q,R)$ & $d\eta$ gives $\operatorname{div}F$ \\
$\int_{\partial M}\omega=\int_{M}d\omega$ & & Green, Gauss, Stokes, FTC \\
\hline
\end{tabular}
\end{center}

One small consequence: because $d(d\omega)=0$, the exterior derivative taken twice is always zero.
Translated, $\operatorname{curl}\nabla f=0$ and
$\operatorname{div}\operatorname{curl}F=0$. A symbolic check for a general function
and a general field gives zero in all six components, so those two identities that look
like coincidences are one statement.

\section{When these theorems fail}
\tingkat{advanced}

All the theorems in this chapter are conditional equalities. This section collects four
ways those equalities break, with numbers.

The first failure is a singularity. Take $F(x,y)=(-y,x)/(x^{2}+y^{2})$
on the plane without the origin. Compute its curl:
\[
\frac{\partial Q}{\partial x}-\frac{\partial P}{\partial y}
=\frac{(x^{2}+y^{2})-2x^{2}}{(x^{2}+y^{2})^{2}}
+\frac{(x^{2}+y^{2})-2y^{2}}{(x^{2}+y^{2})^{2}}=0 .
\]
A finite difference with $h=10^{-5}$ at $(1,0)$ gives $-2.01\times10^{-10}$, zero up to
the truncation error. But the circulation
around the unit circle is $2\pi=6.2831853072$, because there
$F\cdot\gamma'=1$ at every point. So zero curl does not mean zero circulation.

What is violated is the hypothesis of Green: $P$ and $Q$ must have continuous partial
derivatives on an open set containing $D$, whereas both are not
defined at the origin inside $D$, and $D$ is not simply connected. On the annulus $1\le r\le2$ Green still applies, with two boundaries
of opposite direction, and the integral is $2\pi-2\pi=0$; that zero is correct because there
the field has the potential $\theta$, the polar angle.

The second failure is swapping the order without integrability. Take
\[
f(x,y)=\frac{x^{2}-y^{2}}{(x^{2}+y^{2})^{2}}
\]
on the square $(0,1)\times(0,1)$. For fixed $y>0$,
\[
\int_{0}^{1}f(x,y)\,dx=\Bigl[-\frac{x}{x^{2}+y^{2}}\Bigr]_{x=0}^{x=1}
=-\frac{1}{1+y^{2}},
\]
so the outer integral gives $-\pi/4=-0.7853981634$; in the reverse order
the result is $+\pi/4$. Both iterated integrals exist but they differ, so by
Fubini the double integral does not exist. In polar coordinates
$f=\cos2\theta/r^{2}$, and on the unit quarter disk
\[
\iint|f|\,dA\ge\int_{0}^{\pi/2}|\cos2\theta|\,d\theta
\int_{0}^{1}\frac{dr}{r}=\infty ,
\]
because $\int_{0}^{1}dr/r$ diverges at $r=0$. Fubini is not wrong: its hypothesis
is not satisfied.

The third failure is an unmanageable boundary. Gauss demands that $E$ be closed
and bounded with a piecewise smooth boundary. On an unbounded set,
the boundary can run off to infinity and the surface integral stops meaning anything;
Spivak places this as an exercise after Theorem 5-5, which fails without
compactness. Finitely many corners are still safe, because the normal jumps only at
finitely many points; a boundary with endlessly jagged teeth, a coastline for instance, is not.

The fourth failure is a surface that cannot be oriented. On a Möbius
strip, walking around the strip brings the normal back to the opposite
side, so $\mathbf{n}$ in Stokes' formula is not defined.
The theorem is not wrong; its statement is simply incomplete for objects of that kind.

\begin{tantangan}
Take $T(u,v)=(u^{2},v)$ on $D^{*}=(-2,2)\times(0,1)$. This map seems to
widen things: the point $u=2$ lands at $x=4$, and its derivative $2u$ grows without
bound as $u$ moves away from zero. Compute the area of $T(D^{*})$ from the picture, then
compute it with the change-of-variables formula. The two numbers differ by exactly a factor
$2$. Explain why, and name the line on which the Jacobian determinant is zero.

Hint 1: check whether $T$ is one-to-one. There are two values of $u$ that land at
$x=1$.

Hint 2: the line $u=0$ is folded onto a single point, so every image point
is counted twice by the integral; $|\det J|=|2u|$ does not know that the fold
is only one sheet.

If you are stuck, go back to the definition of a valid change of variables and check the
one-to-one condition.
\end{tantangan}

One closing warning, which is not in the textbooks. Numerical quadrature never
checks the hypotheses. It gives $6.2831853072$ with seven correct digits
for a field that violates the hypothesis of Green, and no error message. A tidy number is not proof that the hypothesis holds; it is only proof that
the arithmetic ran.

\section{Worked examples}
\tingkat{intermediate}

Now three computations that form the backbone of this chapter. All three can be checked
with Python, and all three have gone wrong in the author's own hands.

\textbf{One integral, two coordinate systems.}
\nopagebreak[4]

Compute
\[
I=\iint_{D}e^{-(x^{2}+y^{2})}\,dA,\qquad D=\{(x,y):x^{2}+y^{2}\le1\}.
\]
In polar coordinates, $D$ becomes $0\le r\le1$, $0\le\theta\le2\pi$, and
$dA=r\,dr\,d\theta$. The integral separates into two factors:
\[
I=\int_{0}^{2\pi}\!\!\int_{0}^{1}e^{-r^{2}}r\,dr\,d\theta
=2\pi\cdot\Bigl[-\frac12e^{-r^{2}}\Bigr]_{0}^{1}
=2\pi\cdot\frac{1-e^{-1}}{2}=\pi\bigl(1-e^{-1}\bigr).
\]
Its value is $1.9858653038$. The antiderivative with respect to $r$ exists precisely because of the factor
$r$: without that factor, the antiderivative involves the error function.

In Cartesian coordinates, the region demands two pieces:
\[
I=\int_{-1}^{1}\!\!\int_{-\sqrt{1-x^{2}}}^{\sqrt{1-x^{2}}}
e^{-(x^{2}+y^{2})}\,dy\,dx .
\]
The inner integral has no elementary antiderivative, so we have to use
quadrature. The results of the two forms, at the same tolerance:
\begin{align*}
\text{cartesian}&: 1.9858653037988705 \quad (11907\ \text{evaluations}),\\
\text{polar}&: 1.9858653037988712 \quad (441).
\end{align*}
The errors are $8.9\times10^{-16}$ and $2.2\times10^{-16}$; the difference in accuracy
does not matter, but the difference in cost is a factor of $27$.

The midpoint rule on the same uniform Cartesian grid lowers the error more
slowly: $1.141\times10^{-3}$ on $50$ cells per side, and
$4.667\times10^{-5}$ on $400$. Each time the number of cells is doubled, the error
drops by about a factor of $2.9$, not $4$, because the height of the columns near the edge
behaves like $\sqrt{1-x^{2}}$, so the error is of order $h^{3/2}$.

\textbf{Volume of a ball with the Jacobian determinant.}

The volume of a ball of radius $R$ is
\[
V=\int_{0}^{2\pi}\!\!\int_{0}^{\pi}\!\!\int_{0}^{R}
\rho^{2}\sin\varphi\,d\rho\,d\varphi\,d\theta
=2\pi\cdot2\cdot\frac{R^{3}}{3}=\frac{4\pi R^{3}}{3}.
\]
For $R=1$ its value is $4.1887902048$. Three-dimensional quadrature on the ball gives
\[
4.188790204786397 \quad(\text{difference }6.2\times10^{-15}),
\]
and a Monte Carlo simulation with $10^{6}$ points in the cube $[-1,1]^{3}$ gives
$4.192824\pm0.004$.

Compare that with dropping the Jacobian determinant. Without the factor
$\rho^{2}\sin\varphi$,
\[
\int_{0}^{2\pi}\!\!\int_{0}^{\pi}\!\!\int_{0}^{1}1\,d\rho\,d\varphi\,d\theta
=2\pi^{2}=19.7392088022,
\]
which is $3\pi/2$ times too large, and the result is still a number that looks
plausible for a volume.



The following table collects every number in this chapter.

\begin{center}
{\small
\begin{tabular}{llll}
\hline
Quantity & Exact & Numerical & Error \\
\hline
$\iint_{D}e^{-(x^{2}+y^{2})}dA$, polar & $1.985865$ & $1.985865$ & $2.2\times10^{-16}$ \\
$\iint_{D}e^{-(x^{2}+y^{2})}dA$, Cartesian & $1.985865$ & $1.985865$ & $8.9\times10^{-16}$ \\

$\int_{\gamma}F\cdot d\mathbf{r}$ & $10$ & $10.000000$ & $1.8\times10^{-15}$ \\
$\oint_{\partial D}$ Green, square & $2$ & $2.000000$ & $<10^{-15}$ \\
$\iint_{S}(\operatorname{curl}F)\cdot\mathbf{n}\,dA$ & $3.141593$ & $3.141593$ & $<10^{-13}$ \\
\hline
\end{tabular}}
\end{center}

\kertas{
Test yourself in three steps. (1) Recompute $I$ by hand using the
antiderivative $-\tfrac12e^{-r^{2}}$, then compare with $1.9858653038$;
a difference in the last digit is a rounding error, not a formula error. (2)
Run \texttt{scipy.integrate.dblquad} on both forms and count how many
times the function is called; the numbers $11907$ and $441$ will reappear, provided
the tolerance is the same for both. (3) Change $D$ to the disk of radius $2$
and check that the exact answer becomes $\pi(1-e^{-4})$, namely $3.0841$.
If the third step gives a different number, what is wrong is almost certainly the limit on $r$,
not the formula. }

\section{Computational lab}
\tingkat{advanced}

The following code checks five things at once: the cost of the two coordinate systems, Green,
Gauss, Stokes, and a field with a singularity. The numbers in the previous table come
from this program.

\begin{lab}
\begin{verbatim}
import numpy as np
from scipy import integrate

# 1. One double integral, two coordinate systems
def fc(y, x):
    return np.exp(-(x*x + y*y))

cnt = [0]
def fc_count(y, x):
    cnt[0] += 1
    return fc(y, x)

cart = integrate.dblquad(fc_count, -1, 1,
                         lambda x: -np.sqrt(1-x*x),
                         lambda x: np.sqrt(1-x*x),
                         epsabs=1e-12, epsrel=1e-12)
ncart = cnt[0]
cnt[0] = 0

def fp(r, th):
    cnt[0] += 1
    return np.exp(-r*r)*r

pol = integrate.dblquad(fp, 0, 2*np.pi, 0, 1,
                        epsabs=1e-12, epsrel=1e-12)
value = np.pi*(1 - np.exp(-1))
print("cartesian %.16f error %.2e eval %d"
      % (cart[0], abs(cart[0]-value), ncart))
print("polar     %.16f error %.2e eval %d"
      % (pol[0], abs(pol[0]-value), cnt[0]))
print("exact     %.16f" % value)

# 2. Green: (x^2-y)dx + (x+2y)dy on the unit square
P = lambda x, y: x*x - y
Q = lambda x, y: x + 2*y
edges = [integrate.quad(lambda t: P(t, 0.0), 0, 1)[0],
         integrate.quad(lambda t: Q(1.0, t), 0, 1)[0],
         integrate.quad(lambda t: P(t, 1.0), 1, 0)[0],
         integrate.quad(lambda t: Q(0.0, t), 1, 0)[0]]
dbl = integrate.dblquad(lambda y, x: 2.0, 0, 1, 0, 1)[0]
print("Green edges %.15f  double %.15f" % (sum(edges), dbl))

# 3. Gauss: F=(x^3,y^3,z^3) on the unit ball
def flux(phi, th):
    sp, cp = np.sin(phi), np.cos(phi)
    x, y, z = sp*np.cos(th), sp*np.sin(th), cp
    return (x**4 + y**4 + z**4)*sp

F = integrate.dblquad(flux, 0, 2*np.pi, 0, np.pi,
                      epsabs=1e-11, epsrel=1e-11)[0]
print("Gauss flux %.13f  12pi/5 %.13f" % (F, 12*np.pi/5))

# 4. Stokes: F=(z,x,y) on the upper hemisphere of the unit ball
def curl_n(phi, th):
    sp, cp = np.sin(phi), np.cos(phi)
    x, y, z = sp*np.cos(th), sp*np.sin(th), cp
    return (x + y + z)*sp

S = integrate.dblquad(curl_n, 0, 2*np.pi, 0, np.pi/2)[0]
L = integrate.quad(lambda t: np.cos(t)**2, 0, 2*np.pi)[0]
print("Stokes surface %.13f  boundary %.13f" % (S, L))

# 5. A field with a singularity: F=(-y,x)/(x^2+y^2)
def circle(t):
    x, y = np.cos(t), np.sin(t)
    return (-y*(-np.sin(t)) + x*np.cos(t))/(x*x + y*y)

def Pg(x, y): return -y/(x*x + y*y)
def Qg(x, y): return x/(x*x + y*y)

h = 1e-5
curl_01 = ((Qg(1+h, 0) - Qg(1-h, 0))/(2*h)
           - (Pg(1, h) - Pg(1, -h))/(2*h))
C = integrate.quad(circle, 0, 2*np.pi)[0]
print("circle %.13f  2pi %.13f" % (C, 2*np.pi))
print("curl(1,0) %.2e" % curl_01)
\end{verbatim}
\end{lab}

The actual output, as it came out:

\begin{verbatim}
cartesian 1.9858653037988705 error 8.88e-16 eval 11907
polar     1.9858653037988712 error 2.22e-16 eval 441
exact     1.9858653037988714
Green edges 2.000000000000000  double 2.000000000000000
Gauss flux 7.5398223686155  12pi/5 7.5398223686155
Stokes surface 3.1415926535898  boundary 3.1415926535898
circle 6.2831853071796  2pi 6.2831853071796
curl(1,0) -2.01e-10
\end{verbatim}

Four things deserve a reading. First, the two opening lines agree to the fifteenth digit at a cost that differs
by a factor of $27$: the right coordinates do not change the answer, only
its price. Second, Green matches to the last digit, but that is
only proof that the implementation is consistent with its formula. Third, Gauss and Stokes are
computed from different sides and still meet; that is where verification has teeth.
Fourth, quadrature gives $6.2831853071796$ for a field that violates the hypothesis of
Green, without a single error message.

\begin{untukdosen}
The teaching order I use: the density of the plate and the program that forgets the factor
$r$, then Fubini, then the Jacobian determinant as the area of a stretching cell,
then line integrals and their fundamental theorem, and only then Green. Green is best proved
in full for type I regions at the board, because that proof uses only two tools
the students already have. Gauss and Stokes are enough to prove for a box and for a
parametrized surface.

Two misconceptions come up most often. First, students think
$\partial P/\partial y=\partial Q/\partial x$ is enough for conservative, because
they have only tested fields on a domain without holes. The field
$(-y,x)/(x^{2}+y^{2})$ has to appear at the board, not on the reading list. Second,
they treat the Jacobian determinant as a notational formality and write
$r$ or $\rho^{2}\sin\varphi$ as something to memorize without checking the limits.
Compare $4\pi/3$ with $2\pi^{2}$; that ratio of $4.712$ times is more impressive
than a hundred warnings.

One classroom activity: four groups, each computing one side of
Green, Gauss, and Stokes for the same field, then matching numbers. The group
with the sign error is found out within five minutes, and that is better than being told.
The sign of readiness for what comes next: students can explain why orientation
changes the sign without looking at the formula.
\end{untukdosen}

\section{Summary}
\begin{ringkasan}
\begin{itemize}
\item The double integral is defined through upper and lower sums over cells;
Fubini allows the order of integration to be swapped when the function is integrable.
\item Changing coordinates multiplies the area of a cell by $|\det J|$; polar gives
$r$, spherical gives $\rho^{2}\sin\varphi$. Without that factor the result is tidy but
wrong.
\item For a conservative field, the line integral depends only on the endpoints.
\item Green, Gauss, and Stokes state one thing: the integral of a derivative over a region
equals the integral of the field over its boundary, and orientation decides the sign.
\item All those theorems are conditional: singularities, unmanageable boundaries, surfaces that cannot
be oriented, and non-integrable functions make them fail without a warning.
\end{itemize}
\end{ringkasan}

\section{Exercises}
\begin{latihan}
\begin{enumerate}
\item Compute $\iint_{D}(x^{2}+y^{2})\,dA$ on the unit disk in Cartesian and polar coordinates, check it with \texttt{dblquad}, and state the number of evaluations for each method.

\item Prove that if $T(u,v)=A(u,v)$ is linear and $\det A\ne0$, the
change-of-variables formula becomes $|\det A|\iint_{D}f\circ T\,dA$.

\item Repeat the Green example on the square in the clockwise direction: show that the result is
$-2$ and state which side contributes which sign.

\item Compute the volume of the ellipsoid $x^{2}/a^{2}+y^{2}/b^{2}+z^{2}/c^{2}\le1$ with
the change of variables $x=au$, $y=bv$, $z=cw$, then check the formula on
$a=b=c=1$.

\item Assess this claim: true, false, or not enough evidence. ``Because
$\operatorname{curl}F=0$ at every point where $F$ is defined, the field
$F=(-y,x)/(x^{2}+y^{2})$ is conservative on
$\mathbb{R}^{2}\setminus\{(0,0)\}$.'' Include a reason and one supporting
number.

\item Prove that the line integral is independent of path if and only if the
circulation of every closed curve is zero, by using two paths from $A$ to
$B$ that together form one closed curve.

\item Compute both sides of Stokes for $F=(-y,x,z)$ on the paraboloid
$z=1-x^{2}-y^{2}$, $z\ge0$, with the normal pointing upward, then check
with quadrature.

\item Compute both sides of Gauss for $F=(x,y,z)/(x^{2}+y^{2}+z^{2})^{3/2}$ on the unit
ball. Explain why the result does not contradict the theorem, and
which hypothesis is violated. (Hint: compute $\operatorname{div}F$ first.)

\item Show that the Jacobian determinant of spherical coordinates is zero at $\varphi=0$ and
$\varphi=\pi$, then explain why the map may still be used on
$0<\varphi<\pi$. (Hint: check where the map stops being one-to-one.)
\end{enumerate}
\end{latihan}

\begin{untukpeneliti}
What is really proved in this chapter, with a written proof: the lemma on a region
between two graphs, the fundamental theorem for line integrals, its converse through the integral
potential, Green for type I and type II regions, the consequence for the curl test on a
simply connected domain, and the sketch of the proof of Gauss through boxes and of Stokes
through a parametrization. What I quote without proof: Fubini (Spivak, the chapter
\emph{Integration}, Theorem 3-10), the full change of variables (Spivak's same
chapter; Rudin, the chapter \emph{Integration of Differential Forms}), Gauss (Spivak,
Theorem 5-8), and Stokes (Spivak, Theorem 5-9); general Stokes is Theorem
4-13 and 5-5.

What is overpopularized is the phrase ``zero curl means
conservative''. What is correct: zero curl means closed form, and a closed form is
exact only when the first cohomology group of the domain is zero. For $\mathbb{R}^{3}$
that holds, for $\mathbb{R}^{2}\setminus\{0\}$ it does not, and the difference is
measured by the integral on the circle, namely $2\pi$ here.

The open question that is genuinely open for me: can quadrature give
a certificate that the hypotheses of a theorem are satisfied, and not merely an estimate of the value?
The output in the lab shows that a numerical answer correct to fifteen
digits can come from a field that violates the hypotheses, with no warning sign. For
a field that is known only through samples on a grid, there is no general way
to be sure that there is no singularity between the sample points.

Limits of validity: all functions here are real-valued and piecewise
smooth; all numbers come from Python 3.14.4, NumPy 2.5.2, SciPy 1.18.0
with a tolerance of $10^{-12}$ unless stated otherwise, and spherical surfaces
are parametrized with $\mathbf{n}$ pointing outward.
\end{untukpeneliti}

\peta{Chapter 15 follows Stewart, \emph{Calculus: Early Transcendentals}, the chapter ``Multiple Integrals'' (15.1--15.3 and 15.6--15.9) and the chapter ``Vector Calculus'' (16.2
line integrals, 16.3 the fundamental theorem, 16.4 Green, 16.5 curl
and divergence, 16.8 Stokes, 16.9 Gauss); on Spivak, \emph{Calculus on
Manifolds}, the chapter ``Integration'' (Fubini, Theorem 3-10, change of variables),
the chapter ``Integration on Chains'' and the chapter ``Integration of Manifolds'' (Theorems
5-5, 5-8, 5-9); on Rudin, \emph{Principles of Mathematical Analysis}, the chapters
``Functions of Several Variables'' and ``Integration of Differential Forms''; on Active Calculus, chapters 11 and 12 for a more gradual order.
The basis is the chapter \emph{Integral, Series of Functions, and Exchange of Limits} (the integral part)
and the chapter \emph{Continuity and Derivatives} (partial derivatives, Jacobian) of the first edition
of this book. This chapter uses the Riemann integral and the fundamental theorem of calculus
from Chapter 13, uniform continuity and compact sets from Chapter 11, determinants
from Chapter 7, and gradients and Jacobian matrices from Chapter 12, all of them referred to
and not repeated.}
